% Pursuing Stacks (cotes 134-2 to 134-8): archive cote 134-5, site page 48 (PDF page 49), typescript page 235.
% typescript. Draft compared: ps4.tex:2363-2433.
% AI-transcribed by Claude Opus 5.5 (claude-opus-5-5, adaptive thinking, effort high), Message Batches API, batch ps2 (msgbatch_01NVwr1Yuo8qJf5rykeGBXPK), 2026-11-12; not yet checked by eye.
\scanpage{134-5/48}
\mspage{235}

admitting W as it's set of weak equivalences.  We assume beforehand M ad-%
mits a final object $e_M$ .  A ns condition is the \unsure{existence} of a small cate-%
gory A and a functor
\[ i: A \longrightarrow M \]
satisfying the following two conditions:

(i) for x in M of the form i(a) ($a \in \mathrm{Ob}\, A$) \uline{or} $e_M$ , $i^{*}(x) \in A^{\wedge}_{as}$ ,

(ii) $W = (i^{*})^{-1}(\underline{W}_A)$ .

Another ns condition, which looks more pleasant I guess, is that there
exist a small full subcategory B of M, containing $e_M$ , with inclusion
functor $j: B \to M$, such that
\[ (5) \qquad W = (j^{*})^{-1}(\xout{M}\ \underline{W}_B) \quad . \]
In any case, if we want a ns condition on W for it to be the set of
weak equivalences for some asphericity structure on M (not necessarily
an aspheric one, and therefore maybe not unique), we get: there should
exist a small full subcategory B of M, with inclusion functor j, such
that (5) above holds. (Here we do not assume $e_M$ to exist, still less B
to contain it) Similarly, the corollary of prop.1 implies that a sub-%
set $M_{as}$ of Ob M is an asphericity structure on M, iff there exists a
small full subcategory B in M, such that
\[ (6) \qquad M_{as} = (j^{*})^{-1}(B^{\wedge}_{as}) \quad ; \]
moreover, if M admits a final object $e_M$ , the asphericity structure $M_{as}$
is aspheric iff B can be chosen to contain $e_M$ .

\uline{Remarks} 2.  The relationship between aspheric asphericity struct\xout{u}ures
$M_{as}$ on M, and sets $W \subset F\!\!/(M)$ of ``weak equivalences'' in M \xout{(}satisfying
the condition above (and which we may call ``weak equivalence structures''
on M), in case M admits a final object $e_M$ , is reminiscent of the rela-%
tionship between \xout{thomste} ``contractibility structures'' $M_{c\xout{ont}} \subset \mathrm{Ob}\, M$ on M,
and those \xout{h}``homotopism structures'' \xout{xxxM} $h_M \subset F\!\!/(M)$ on M which can be
described in terms of such a contractibility structure (cf. sections 51
and 52). Both pairs can be viewed as giving two equivalent ways of
expressing one and the same kind of structure - the structure concerned
by the first pair being centered on asphericity notions, whereas the
second is concerned with typical homotopy notions rather. We'll see
later that any contractibility structure defines in an evident way an
\margin{x} asphericity structure, and in the most interesting cases (e.g. canonical
\xout{contractibl} modelizers), it is uniquely determined by the latter.

\margin{48}
